\documentclass[11pt,letterpaper]{article}
\usepackage{amssymb}

\begin{document}

\begin{center}
{\Large \bf Numerical Linear Algebra }\\
{\Large \bf Weekly assignment \#2 }
\end{center}

\small

\begin{enumerate}

\item
Write a Julia code that starts with something like
\begin{verbatim}
using Random
using LinearAlgebra
m = 10; n = 7;
A = randn(rng,Float64,m,n);
b = randn(rng,Float64,m,1);
Q, R = qr(A)
\end{verbatim}
The goal is to find $x$ in $\mathbb{R}^n$ such that $x$ realizes the minimum of 
$\min_{x\in\mathbb{R}^n} \| b - Ax \|_2$ using the computed 
factors $Q$ and $R$. Check that the answer $x$ is correct.
Note that $Q$ might never be explicitly computed and you can manipulate it
with $Q'*x$ for example, and that will return what you expect.

\item
Write a Julia code that starts with something like
\begin{verbatim}
using Random
using LinearAlgebra
rng = MersenneTwister()
m = 7; n = 10;
A = randn(rng,Float64,m,n);
b = randn(rng,Float64,m,1);
\end{verbatim}
The goal is to find $x$ in $\mathbb{R}^n$
such that $x$ is the $x$ of minimum 2-norm such that $Ax=b$.
using a QR factorizarion and the computed 
factors $Q$ and $R$. 
In other words, $x$ realizes
$\min_{x\in\mathbb{R}^n, Ax=b} \| x \|_2$.
Check that the answer $x$ is correct.

\item
Write a Julia code that starts with something like
\begin{verbatim}
using Random
using LinearAlgebra
m = 10; n = 7; k = 4;
A = randn(rng,Float64,m,k) * randn(rng,Float64,k,n);
b = randn(rng,Float64,m,1);
U, S, V = svd(A)
\end{verbatim}
Note that $A$ is constructed such that the rank of $A$ is $k=4$.
The goal is to find $x$ in $\mathbb{R}^n$ such that $x$ 
is the vector of minimum norm that 
realizes the minimum of 
$\min_{x\in\mathbb{R}^n} \| b - Ax \|_2$
using the computed 
factors $U$, $S$ and $V$.
Check that the answer $x$ is correct.\\
\underline{Note:} $x$ can also be computed with a 
``complete orthogonal decomposition''. 

\item 
Here is an algorithm to compute a QR factorization of $A$.
\begin{verbatim}
using Random
using LinearAlgebra
rng = MersenneTwister()
m = 10;
n = 7;
A = randn(rng,Float64,m,n);
Q = copy(A)
R = randn(rng,Float64,n,n);
BLAS.syrk!( 'U', 'T', 1.0, Q, 0., R )
LAPACK.potrf!( 'U', R )
BLAS.trsm!( 'R', 'U', 'N', 'N', 1.0, R, Q )
\end{verbatim}
Check that this algorithm works. (I.e. check that the computed $Q$ and $R$ factors
are QR factors of a QR factorization of $A$.)
Explain how this three-line algorithm works and why it works. 
(By typing your explanations in
your file as some "comments".) Compute the number of operations of this
algorithm as a function of $m$ and $n$. Change the condition number of the 
input matrix $A$ and observe how $\|I-Q^TQ\|_2$ depends on the condition number.
What do you think is the relation of $\|I-Q^TQ\|_2$ as a function of the condition number?
Time this code for a matrix $m=5,000$ and n=$1,000$. Compute the performance in GFLOP/sec.
Compare with built-in julia QR. (You can leave $Q$ in the ``Householder vector'' form.)
Or you can call LAPACK directly with
\begin{verbatim}
tau = randn(rng,Float64,n);
LAPACK.geqrf!( A, tau )
\end{verbatim}

\item We start with
\begin{verbatim}
m = 10; n = 7; A = randn(rng,Float64,m,n); 
u = randn(rng,Float64,m,1);
v = randn(rng,Float64,n,1);
Q, R = qr(A)
Q = Matrix(Q)
\end{verbatim}
We now consider the matrix $B=A+uv^T$. $B$
is a rank-1 update of $A$. Design an algorithm to find the QR factorization of $B$ in
$\mathcal{O}(n^2)$ operations by updating the $Q$ and $R$ factors of $A$.
Implement your algorithm in Julia and verify the correctness.

\end{enumerate}






\end{document}
